\section{Empirical Results}\label{sec:results}

\subsection{Results on the Independent Panel}\label{sec:independent_results}

Table~\ref{tab:independent_panel} presents the 432 observations from the prospective comparison. All six Holm-adjusted loss comparisons have $p>0.05$. In inventory all three mean losses are within 0.0024. Corrected KG uses 369.0 target calls on average and certifies all 48 runs. PublicPsi uses 396.4 calls and certifies 47. The paired reduction is 27.42 calls per run (individual 95\% interval $[5.61,49.22]$). Each arm uses 136 search calls per run, so the inventory savings arise during verification. KG takes 15.52 seconds per run on average, compared with 7.97 for PublicPsi.

In queueing, corrected KG uses 352.3 mean target calls, compared with 365.6 for GP-FTA-TS and 374.6 for PublicPsi. The 22.33-call reduction relative to PublicPsi has an individual 95\% paired interval of $[-0.52,45.18]$. KG wins 19 of 24 paired task-mean loss comparisons against PublicPsi, with a mean difference of 0.0325 (individual 95\% interval $[0.0047,0.0603]$, Holm-adjusted $p=0.119$). The mean loss of the same-framework TS arm is within 0.0029 of KG. All 144 deployments are certified and classified as confidently safe.

In traffic, PublicPsi has lower mean loss than KG by 0.0455 (individual 95\% interval $[0.0101,0.0809]$, Holm-adjusted $p=0.084$), uses 23.83 fewer calls on average (individual 95\% paired interval $[-4.11,51.78]$), and certifies 42 of 48 runs. All 42 are confidently safe. KG certifies 32, of which 29 are confidently safe, two boundary-ambiguous and one confidently unsafe. The illustrative 10\% abstention surcharge preserves this traffic ordering and places KG first by mean loss in inventory.

The 72 saved status-quo audits used 2,000 replications per task and reproduce the status-quo costs in all 432 run records. All 24 anchors in each of inventory and queueing are confidently safe. All 24 stylized traffic anchors are confidently unsafe, with violation point estimates from 22.80\% to 95.95\%. The lowest traffic estimate is 456 violations in 2,000 audited replications (95\% interval $[20.98\%,24.70\%]$). When the operator retains the anchor after abstention, the resulting traffic policy is confidently unsafe in 17 of 48 KG runs, 11 of 48 GP-FTA-TS runs, and six of 48 PublicPsi runs (Table~\ref{tab:operating_policy}). For KG the 17 comprise one certified nominee and 16 retained anchors. Those 16 anchors have point estimates from 50.25\% to 95.95\%. These traffic results illustrate how certification of new policies and safety of the policy actually operated can diverge.

Recomputing the exact test and choice rule on 885 saved verification slots reproduces the recorded terminal decisions and calls. The rule deploys a non-primary nominee in 51 of the 128 corrected KG certifications, including 20 of 32 in traffic. Of these 51 deployments, 30 selected a later slot despite a primary pass, and 21 followed a primary failure. The final decision thus frequently uses a policy other than the first nominee scored by the search proxy.

The recorded search counters classify a query as a revisit when its normalized policy, rounded to six decimal places, matches a policy already observed. Among the 2,880 online queries per method and family, the inventory, queueing, and traffic revisit rates are 33.2\%, 30.1\%, and 37.4\% for KG. The corresponding rates are 39.6\%, 37.2\%, and 40.6\% for GP-FTA-TS, and 41.9\%, 39.2\%, and 42.7\% for PublicPsi. These rates quantify the use of adaptive pooling. They leave the size of any resulting finite-sample variance bias open.

The panel used 165,864 target calls, comprising 58,752 for search and 107,112 for terminal verification. The independent audits used 1,320,000 calls. Source archives and fitted priors were reused. This comparison measures the corrected implementation on newly drawn tasks at the common cap. It does not estimate the causal effect of the quadrature adjustment in isolation.

\subsection{Matched Results on Queueing and Traffic}\label{sec:matched_results}

Tables~\ref{tab:matched_policy} and~\ref{tab:matched_contrasts} report the four paid arms and the exact replay of the early release policy. All ten Holm-adjusted $p$ values equal 1. In queueing, the historical diagonal prior raises the mean loss by $0.0380$ relative to the target-initialized weak prior. Adding off-diagonal covariance lowers mean loss by $0.0222$ (individual 95\% interval $[-0.0483,0.0926]$ for the reduction). In traffic, full covariance raises mean loss relative to diagonal by $0.0147$ (reduction interval $[-0.0393,0.0099]$). The point estimates show how the contribution of the historical coefficients can change with the prior and problem.

Full covariance changes the first query in 7 of 48 queueing runs and 18 of 48 traffic runs. Relative to the diagonal prior, it changes 46 primary nominations in each family and changes the operating policy in 46 queueing and 35 traffic runs. In queueing all policies operated by both arms are confidently safe. Full covariance lowers loss in 25 of those paired runs and raises it in 21. In traffic it turns three confidently unsafe operating policies into confidently safe ones, and six confidently safe policies into confidently unsafe ones. When both traffic runs are confidently safe, full covariance lowers loss in 12 and raises it in 13. All 24 traffic anchors on this new panel are confidently unsafe, and all confidently unsafe operating outcomes in Table~\ref{tab:matched_policy} arise from retaining an anchor after abstention.

Keeping the full coefficient model fixed, weighting by gate power rather than chance feasibility raises mean loss by $0.0314$ in queueing and $0.0208$ in traffic. Mean terminal use rises by 1.0 call in queueing and falls by 7.5 in traffic. Queueing yields 48 confidently safe policies with either weighting. Traffic yields 30 confidently safe operating policies with gate power and 33 with feasibility weighting. These comparisons assess two scores used with the same terminal test.

The initial score releases the historical coefficients in five queueing runs and three traffic runs out of 48 in each family. Relative to always using the full prior, the replayed strategy reduces mean queueing loss by $0.00003$ and raises mean traffic loss by $0.00106$. Its mean loss is higher than always using the weak coefficients in both families. Of the three traffic releases, using weak coefficients raises loss twice and ties once. Of the 45 traffic runs that retain the full prior, the paired weak arm has lower loss in 25. The predictive score evaluates the constraint response in the common initialization. The operating loss also depends on objective comparisons, candidate nomination and the cost of abstention. After averaging streams within each task, the descriptive Spearman correlations of the initial weak-minus-full log score with the full-minus-weak operating loss are $-0.06$ in queueing and $-0.22$ in traffic. We retain the registered threshold.

As an exploratory direct comparison under the common basis, full historical minus weak-prior operating loss is $0.0158$ in queueing (individual 95\% interval $[-0.0391,0.0707]$) and $0.0274$ in traffic ($[-0.0114,0.0661]$). The earlier queueing comparison against PublicPsi varies several other components and does not identify the benefit of coefficient coupling.

Across the paid arms, online revisit rates range from $31.5\%$ to $35.0\%$ in queueing and $38.2\%$ to $42.4\%$ in traffic. Comparing variance estimates to the independent audit at nominated policies, the fixed-batch constraint estimate has a lower mean squared error in traffic and a higher one in queueing. For certified deployments the pooled estimate has lower mean squared error in both families. Electronic Companion~\ref{ec:matched_variances} reports the corresponding values and sampling scope. The matched study uses 52,224 search, 97,344 verification, and 1,208,000 audit calls across the four paid arms and their shared anchor audits. It reuses the fitted sources.

\subsection{First-slot and full verification on the same search paths}\label{sec:terminal_replay}

Table~\ref{tab:terminal_replay} recomputes operating loss from all 384 saved runs and their exact gate outcomes. The first-slot counterfactual retains the realized search, primary nomination, 112-replication test and original error allocation. It operates the primary on a pass and the anchor on a failure. The actual rule may deploy a later candidate. In traffic, gate-power and feasibility KG have first-slot losses of $1.17594$ and $1.18329$, with 23 and 21 passes. Full verification increases certification counts to 31 and 33, and their mean losses become $1.13071$ and $1.10987$. The gate-power minus feasibility loss difference changes from $-0.00735$ (individual 95\% interval $[-0.04204,0.02735]$) to $0.02084$ ($[-0.00897,0.05066]$). Both first-slot outcomes use 112 terminal calls. The full rule averages 276.75 and 284.25 calls, respectively. The ordering of the sample means reverses. In queueing the feasibility-weighted arm has lower mean loss under both rules. This replay describes the effect of changing the terminal rule on already realized paths. Both acquisitions use the value of a single candidate during search.

In traffic, the full-minus-diagonal mean loss difference is $0.01468$. The ten paired streams in which one arm certifies and the other abstains contribute $0.01440$ to that difference, with contributions divided by all 48 pairs. Both arms certify in 27 pairs and both abstain in 11. For gate-power versus feasibility weighting, the six pairs with different certification outcomes contribute $0.01706$ of the $0.02084$ mean difference. In this sample, the disagreement about which policies to certify often determines whether the operator retains the unsafe anchor. These are arithmetic contributions to the sample averages.

\subsection{Source Sensitivity}\label{sec:two_mechanisms_results}

Refitting on all six source tasks reproduces the archived fits' predictions at the fixed policies. Deleting one source task can change the predicted cost or constraint ordering substantially. Among the 46 distinct full-versus-diagonal primary pairs in queueing, at least one deletion reverses the prior mean cost order in 38 and the constraint order in 35. Traffic yields 45 and 44 reversals out of 46. The median initial absolute cost gaps are $0.0980$ and $0.0280$. Across deletions the median absolute changes in those gaps range from $0.0571$ to $0.2594$ in queueing and from $0.0317$ to $0.1325$ in traffic. After conditioning on a common hypothetical observation one predictive standard deviation above the original status-quo mean, the effect of using full covariance on the cost gap retains its direction across all deletions in two queueing pairs and no traffic pairs. The analogous constraint-gap counts are four and zero. These sensitivity diagnostics identify dependence on particular source tasks, including when the predicted association is propagated to comparisons between target policies.

After conditioning on the eight restored target batches, at least one of six deletions reverses $700/844$ full-covariance cost orders in queueing and $809/857$ in traffic. The corresponding diagonal counts are $708/844$ and $807/857$. Averaging the reversal fractions across the six separate deletions gives $30.0\%$ and $45.2\%$ for full covariance and $31.7\%$ and $46.7\%$ for diagonal covariance. Full-covariance constraint orders reverse in at least one fold for $650/844$ and $787/857$ pairs. These descriptive counts compare shared fits and policies.

Within the original six-source fit, adding only the non-diagonal terms reverses 19 cost orders and two constraint orders in queueing, and 77 and 53 in traffic. None of these reversals preserves both opposing orders in all six deletions. Taken together, the comparisons show sensitivity of the entire fitted prior and instability of decisions specifically changed by the non-diagonal terms. A rule that uses deletion stability to decide when to use those terms would require separate validation.

\subsection{Scope of the fixed-roster comparison}\label{sec:relaxed_eligibility}

Among the 96 earlier fixed rosters, the primary and third-slot candidates were unevaluated in every case and distinct within every roster. The second slot was already evaluated in every case. The five- and six-decimal keys agree on this classification. Since one more query evaluates only one policy, none of the rosters can be nominated immediately after a query under the executed protocol. The bounds in Table~\ref{tab:query_screen} apply to those four frozen rosters per state as a relaxed approximation to a later terminal menu.

The original query and the refined first-slot choice on those menus disagree in $9/12$ queueing and $5/12$ traffic states. Within the same menus, changing from the first-slot to the full-tree value changes $3/12$ and $6/12$ choices. The original query differs from the refined full-tree choice in $9/12$ and $7/12$ states. The within-menu comparison isolates the valuation of later slots. The fixed-roster diagnostic also permitted immediate abstention and held the fitted noise model fixed in its fantasies. The comparisons with the original query also change the menu and the numerical score.

The relaxed screen identifies a unique query in ten states. In traffic, three states retain one query, one retains two and eight retain 155 to 160 out of 160. One traffic state accounts for $78.45\%$ of the sample's estimated gain from changing choices. The detailed timing and the refinement after coarse quadrature appear in Electronic Companion~\ref{ec:query_diagnostics}. The observed backup value and audit risk do not by themselves supply a threshold for switching scores.

\subsection{One more query under the actual nomination rule}\label{sec:eligible_queries}

Under the designated 64-node forecast, the two scores choose different queries in one of the 12 queueing states and none of the 12 traffic states at initialization. The screen uniquely determines the full-rule choice in all 12 queueing states and 11 of 12 traffic states. Its only initial switch is queueing task 8. The first-slot score chooses policy index 24 and the full-rule score chooses 124. Keeping the first-slot choice gives up between $0.001044$ and $0.004474$ modeled value on this finite law, with a numerical estimate of $0.002849$. In the one nonunique case, traffic task 0, both numerical scores choose index 76. The screen retains 140 of 160 candidates, and the value given up by retaining the first-slot choice is at most $0.000443$. The results show how either a unique winner or a small regret bound can guide computation.

All 24 primary-law first-slot winners agree across the two marginal integration orders, and their full-rule winners agree across both sizes and both scrambles of the terminal integral. The disagreement status agrees across all four outer rules in 11 of 12 states in each family. Across the two outer scrambles and two resolutions the disagreement count is between zero and one per family. The interval statements are conditional on their respective finite forecasts. Median CPU time to update the model and nominate under both sets of fantasies is $68.46$ seconds in queueing and $71.64$ in traffic. The marginal calculation and screening medians are $2.41$ and $1.98$ seconds, while the joint-integration medians are $0.20$ and $0.17$ seconds. These times measure the entire diagnostic across 160 queries per state. In this sample the dominant work is updating the model and forming the legal shortlist at every fantasy.

In the four near-terminal histories, the two scores choose different queries in one of two queueing states and both traffic states (Table~\ref{tab:eligible_query}). Their estimated gains from selecting the full-rule query are $0.002112$ in the switching queueing state and $0.000022$ and $0.001760$ in the traffic states. Every full-rule winner agrees across all registered outer and inner comparisons in these four histories. The three disagreements persist across outer forecasts, and one additional disagreement appears under the second scramble. The interval screen retains $160,160,160,159$ candidates respectively and determines no unique choice. Median joint integration rises to $18.42$ and $15.11$ seconds in queueing and traffic, with end-to-end medians of $160.28$ and $129.07$ seconds. These four paths describe behavior near deployment and do not estimate its frequency in the population.

The earlier relaxed rosters give more initial disagreements under their working forecast. The new initial results share the same proposal queries but also use observation-dependent nomination, signed savings, and HVD refitting after both repeats. The change between studies is therefore not a pure effect of removing unevaluated nominations. Within the new design, the first-slot and full-rule scores share all these features. Their differences isolate the modeled effect of valuing later slots under the actual eligibility constraint.

\subsection{Earlier shared-panel evidence}

The frozen U3 runtime has lower mean loss than each of six raw-policy comparators in four families on the previously used P2 panel at $N=60$. Comparisons with the same-representation and same-framework controls are mixed. Electronic Companion~\ref{ec:shared_panel} preserves these development and shared-panel results, including the safety-law checks. Archive heterogeneity, microscopic traffic and amortization results appear in Electronic Companion~\ref{ec:archived_sensitivity}.
